\documentclass[11pt]{article}
\usepackage[review]{acl}
\usepackage{times}
\usepackage{latexsym}
\usepackage[T1]{fontenc}
\usepackage[utf8]{inputenc}
\usepackage{microtype}
\usepackage{inconsolata}
\usepackage{graphicx}
\usepackage{amsmath}
\usepackage{amssymb}
\usepackage{booktabs}
\usepackage{multirow}
\usepackage{xcolor}
\usepackage{url}

\title{Why Lightweight Gating Cannot Fix Multi-Turn Attention Degradation:\\ A Mechanistic Analysis}

\author{Anonymous}

\begin{document}
\maketitle

\begin{abstract}
Large Language Models (LLMs) progressively lose attention to goal-relevant tokens as multi-turn dialogues grow longer, a phenomenon quantified by the Goal Attention Ratio \citep[GAR;][]{dongre2026}. Prior work shows that goal information persists in the residual stream (probe AUC up to 0.98) even as the attention channel closes, motivating a natural intervention: learnable per-layer gates on the residual connection trained with an auxiliary GAR objective. We implement this hypothesis as ThreadGate, a lightweight residual gating module jointly optimized with next-token prediction (NTP) and GAR losses. ThreadGate successfully raises GAR from 0.02 to 0.48 during training, yet this improvement completely fails to transfer to inference (GAR = 0.016), revealing a 30$\times$ training--inference gap. We identify two coupled mechanisms behind this failure: (1) NTP gradients dominate gate parameters by 7--15$\times$ in magnitude, causing gates to learn \emph{reverse turn conditioning} that suppresses intervention where it is most needed; (2) at inference, gates attenuate sublayer outputs by only 2--7\%, rendering the gating mechanism effectively inert. Four-way ablation across two model families (Qwen-2.5-7B and Llama-3.1-8B) confirms that gating provides zero incremental benefit over LoRA-only fine-tuning (paired bootstrap $p = 0.71$, Cohen's $d = 0.02$). A scheduled sampling experiment reduces the gap from 30$\times$ to 4.3$\times$, providing causal evidence that exposure bias is the primary driver. Our findings constitute the first quantitative measurement of how teacher-forcing distribution shift destroys auxiliary attention loss generalization in lightweight residual gating.
\end{abstract}

%---------------------------------------------------------------------
\section{Introduction}
%---------------------------------------------------------------------

Multi-turn dialogue is the dominant interaction paradigm for LLMs, yet model performance degrades as conversations grow longer. \citet{dongre2026} formalized this degradation through the Goal Attention Ratio (GAR), demonstrating 27--48\% drops across 12 models over 50 turns. Their residual stream probes reveal a striking asymmetry: goal information persists in the hidden states (AUC up to 0.98) even as attention to goal tokens collapses. This asymmetry creates a compelling hypothesis: if goal information survives in the residual stream but attention fails to access it, a learnable gate on the residual connection should recover it.

We test this hypothesis through ThreadGate, a per-layer gating module that modulates sublayer contributions via a ReZero-style mechanism \citep{bachlechner2021rezero}. Trained jointly with NTP and auxiliary GAR losses, ThreadGate achieves its optimization objective convincingly---GAR rises from 0.02 to 0.48 during training. However, at inference time under autoregressive generation, GAR drops back to 0.016, a 30$\times$ gap that renders the entire training-time improvement moot.

This gap is not a bug. We verify through line-by-line code audits and controlled experiments that the discrepancy arises from a fundamental distribution shift: attention patterns optimized under teacher-forced contexts do not transfer to autoregressive generation, where the model conditions on its own outputs. While exposure bias is a known phenomenon in sequence generation \citep{bengio2015scheduled, lamb2016professor}, its impact on \emph{auxiliary attention losses}---which operate on internal attention distributions rather than output sequences---has not been quantified.

We identify two mechanisms that amplify this distribution shift. First, NTP loss gradients dominate GAR loss gradients on gate parameters by 7--15$\times$ in magnitude, with the imbalance growing during training as GAR loss saturates. This causes gates to learn \emph{reverse turn conditioning}: suppressing sublayer outputs at early turns (where GAR is naturally high) while becoming passive at late turns (where intervention is needed). Second, at inference time, gate values cluster near 1.0 (mean 0.93--0.98), attenuating sublayer outputs by only 2--7\%---far too little to redirect attention patterns. Four-way ablation (ThreadGate vs.\ Uniform Gate vs.\ LoRA-only vs.\ Base) demonstrates that gating provides zero incremental benefit: all fine-tuned variants converge to identical GAR and perplexity.

We further validate these findings on Llama-3.1-8B \citep{dubey2024llama3}, confirming that the 26--30$\times$ gap, gradient competition, and zero gating benefit are architecture-independent. A scheduled sampling experiment \citep{bengio2015scheduled} reduces the gap from 30$\times$ to 4.3$\times$ (a 7$\times$ reduction), providing causal evidence that exposure bias accounts for the majority of the training--inference disconnect.

Our contributions are:
\begin{enumerate}
    \item The first quantitative measurement of how teacher-forcing distribution shift destroys auxiliary attention loss generalization in lightweight residual gating (30$\times$ GAR gap).
    \item Per-layer quantification of NTP--GAR gradient competition: 67--79\% of layers exhibit directional conflict, with NTP magnitude dominance growing from 7.5$\times$ to 14.7$\times$ during training.
    \item Systematic ablation across two architectures (Qwen-2.5-7B, Llama-3.1-8B) isolating gating, fine-tuning, and loss design as independent factors, with bootstrap statistical validation ($n=50$, $p=0.71$).
    \item Causal evidence via scheduled sampling that exposure bias is the primary driver of the gap (30$\times \to$ 4.3$\times$), upgrading the result from pure negative finding to diagnosis with partial mitigation.
\end{enumerate}

%---------------------------------------------------------------------
\section{Related Work}
%---------------------------------------------------------------------

\paragraph{Multi-turn dialogue degradation.}
Performance degradation in extended dialogues has been documented across instruction following \citep{zheng2024judging}, factual consistency, and goal adherence \citep{dongre2026}. Related phenomena in long-context settings include the lost-in-the-middle effect \citep{liu2024lost} and attention sink concentration \citep{xiao2024efficient}. \citet{dongre2026} provided the first systematic quantification through GAR and demonstrated that goal information persists in residual streams despite attention collapse, motivating architectural interventions on the residual pathway.

\paragraph{Residual gating and adaptive depth.}
Learnable gates on residual connections have been explored for adaptive computation depth \citep{bachlechner2021rezero, veit2016residual} and layer-wise importance estimation. ReZero \citep{bachlechner2021rezero} initializes gates at zero for stable deep training; SkipNet \citep{wang2018skipnet} learns binary skip decisions. Our ThreadGate adapts this mechanism for attention steering, but we find that the auxiliary loss training signal is too weak relative to NTP to produce meaningful gate modulation at inference time.

\paragraph{Attention steering and intervention.}
Methods for directing transformer attention include attention supervision \citep{liu2016attention}, activation editing \citep{meng2022locating, li2024inference}, and representation engineering \citep{zou2023representation}. These approaches differ in whether they operate at training time (attention losses), inference time (activation patching), or at the representation level. Our work shows that training-time attention losses face a fundamental generalization barrier due to teacher-forcing distribution shift, suggesting that inference-time approaches may be more promising for multi-turn attention recovery.

\paragraph{Exposure bias and scheduled sampling.}
The mismatch between teacher-forced training and autoregressive generation was identified by \citet{williams1989learning} and formalized as exposure bias by \citet{ranzato2016sequence}. Mitigation strategies include scheduled sampling \citep{bengio2015scheduled}, Professor Forcing \citep{lamb2016professor}, and sequence-level training \citep{leblond2018searnn, ranzato2016sequence}. These methods primarily target output sequence quality; we demonstrate that exposure bias also critically affects \emph{internal} auxiliary losses that supervise attention distributions, a setting not previously studied.

%---------------------------------------------------------------------
\section{Background}
\label{sec:background}
%---------------------------------------------------------------------

\paragraph{Goal Attention Ratio (GAR).}
Following \citet{dongre2026}, we define GAR as the fraction of attention mass allocated to goal-relevant tokens in the input. Let $\mathcal{G}$ denote the set of goal token positions (system prompt and initial user query) and $\alpha_{i,j}$ the attention weight from response token $i$ to position $j$. For a set of response tokens $\mathcal{R}$:
\begin{equation}
    \text{GAR} = \frac{1}{|\mathcal{R}|} \sum_{i \in \mathcal{R}} \sum_{j \in \mathcal{G}} \alpha_{i,j}
\end{equation}
We use two instantiations: \textbf{GAR-Dongre} (goal = system prompt task tokens, response = all response tokens) for comparability with prior work, and \textbf{GAR-ours} (goal = system prompt + first user query, response = last token per assistant turn) as the training signal, which provides a richer goal representation and more stable gradient signal.

\paragraph{Residual stream persistence.}
\citet{dongre2026} demonstrated via linear probes that goal-relevant information persists in the residual stream throughout multi-turn conversations (probe AUC up to 0.98), even as attention-based GAR drops by 27--48\%. This decoupling between information availability and attention access motivates our residual gating approach: if the information is present but inaccessible via attention, modulating the residual pathway should enable recovery.

%---------------------------------------------------------------------
\section{Method}
\label{sec:method}
%---------------------------------------------------------------------

\subsection{ThreadGate Architecture}

ThreadGate adds a learnable scalar gate $g_l \in \mathbb{R}$ to each transformer layer $l$, modulating the sublayer contribution via a ReZero-style mechanism \citep{bachlechner2021rezero}:
\begin{equation}
    \mathbf{x}_{l+1} = \mathbf{x}_l + \sigma(g_l) \cdot f_l(\mathbf{x}_l)
\label{eq:gate}
\end{equation}
where $f_l$ denotes the sublayer computation (self-attention or FFN), $\sigma$ is the sigmoid function constraining the gate to $[0, 1]$, and $g_l$ is initialized such that $\sigma(g_l) \approx 1.0$ (preserving the pretrained model's behavior). When $\sigma(g_l) \to 0$, the sublayer output is suppressed and the layer becomes a no-op (identity on the residual stream). When $\sigma(g_l) \to 1$, the layer operates normally. ThreadGate is combined with LoRA adapters \citep{hu2022lora} ($r=16$) on attention projections for parameter-efficient fine-tuning.

\subsection{Loss Design}
\label{sec:loss}

We explored seven loss configurations (v1--v7) spanning four design families: cumulative attention (v1--v2), GAR maximization (v3--v5), smooth saturation (v6), and MSE target (v7). Only v7 achieved stable training convergence.

The v7 loss combines three terms:
\begin{equation}
    \mathcal{L} = \mathcal{L}_\text{NTP} + \lambda_\text{GAR} \cdot \mathcal{L}_\text{GAR} + \lambda_\text{div} \cdot \mathcal{L}_\text{div}
\end{equation}
where $\mathcal{L}_\text{NTP}$ is standard next-token prediction loss, $\mathcal{L}_\text{GAR} = (\text{GAR} - \tau)^2$ is an MSE loss driving GAR toward a target $\tau = 0.5$, and $\mathcal{L}_\text{div}$ encourages diversity across layer-wise gate values to prevent collapse. The MSE formulation provides bidirectional gradients (unlike maximization losses that have unbounded gradients) and a clear convergence criterion.

\subsection{Training Setup}

We fine-tune Qwen-2.5-7B-Instruct \citep{yang2024qwen2} and Llama-3.1-8B-Instruct \citep{dubey2024llama3} on a 1,000-conversation subset of UltraChat \citep{ding2023enhancing} for 3 epochs. LoRA rank is 16, applied to all attention projection matrices. Gate parameters use a separate learning rate ($5 \times 10^{-4}$). All experiments use seed 42.

\subsection{Ablation Design}
To isolate the contribution of gating, we evaluate four configurations:
\begin{itemize}
    \item \textbf{ThreadGate (TG)}: LoRA + per-layer learned gates + NTP + GAR loss (Eq.~\ref{eq:gate}).
    \item \textbf{Uniform Gate (UG)}: LoRA + a single shared gate across all layers + NTP + GAR loss. Isolates whether per-layer specialization matters.
    \item \textbf{LoRA-only (LO)}: LoRA + NTP loss only, no gates. Tests whether fine-tuning alone suffices.
    \item \textbf{Base}: Pretrained model without any fine-tuning.
\end{itemize}

%---------------------------------------------------------------------
\section{Experiments}
\label{sec:experiments}
%---------------------------------------------------------------------

\subsection{GAR Degradation Is Real and Severe (C1)}

We first confirm GAR degradation on our evaluation models. On Qwen-2.5-7B, GAR drops from 0.042 at turn 1 to 0.008 at turn 10 (raw drop 70.8\%, normalized drop 21\%). On Llama-3.1-8B, GAR drops from 0.035 to 0.006 (raw drop 69\%, normalized drop 18\%). Table~\ref{tab:gar-degradation} shows the turn-by-turn comparison. The consistency across architectures (Qwen uses standard MHA; Llama uses GQA with 8 KV heads) confirms that GAR degradation is a structural phenomenon of autoregressive transformers in multi-turn settings, not an architecture-specific artifact.

\begin{table}[t]
\centering
\small
\begin{tabular}{lcccc}
\toprule
\textbf{Turn} & \textbf{Qwen} & \textbf{Llama} \\
\midrule
1 & 0.042 & 0.035 \\
2 & 0.024 & 0.020 \\
3 & 0.021 & 0.017 \\
5 & 0.013 & 0.011 \\
7 & 0.010 & 0.007 \\
10 & 0.008 & 0.006 \\
\midrule
Relative drop & 70.8\% & 69\% \\
\bottomrule
\end{tabular}
\caption{GAR-Dongre degradation across turns for Qwen-2.5-7B and Llama-3.1-8B base models. Both architectures show $>$69\% relative drop.}
\label{tab:gar-degradation}
\end{table}

\subsection{Training Succeeds, Inference Fails (C2)}

ThreadGate v7 converges successfully on both models, exhibiting a characteristic three-phase training dynamic: plateau (steps 0--90), rapid rise (steps 90--170), and convergence. Table~\ref{tab:train-infer-gap} shows the central result: training GAR reaches 0.48 (Qwen) and 0.42 (Llama), but inference GAR collapses to 0.016 on both, yielding 30$\times$ and 26$\times$ gaps respectively. Line-by-line code verification and controlled experiments confirm this is not an implementation artifact.

\begin{table}[t]
\centering
\small
\begin{tabular}{lccc}
\toprule
\textbf{Model} & \textbf{Train GAR} & \textbf{Infer GAR} & \textbf{Gap} \\
\midrule
Qwen-2.5-7B & 0.48 & 0.016 & 30$\times$ \\
Llama-3.1-8B & 0.42 & 0.016 & 26$\times$ \\
\bottomrule
\end{tabular}
\caption{Training--inference GAR gap. GAR-ours is reported for training (teacher-forced context); GAR-Dongre for inference (autoregressive generation).}
\label{tab:train-infer-gap}
\end{table}

\subsection{Gating Provides Zero Incremental Benefit (C4)}
\label{sec:four-way}

Table~\ref{tab:four-way-qwen} presents the four-way ablation on Qwen-2.5-7B. All three fine-tuned variants (TG, UG, LO) converge to nearly identical GAR and perplexity (PPL), while the base model lags substantially. Critically, ThreadGate provides no measurable improvement over LoRA-only fine-tuning on either metric, demonstrating that the gating mechanism is entirely inert at inference time.

\begin{table}[t]
\centering
\small
\begin{tabular}{lcccc}
\toprule
\textbf{Model} & \textbf{GAR-D} & \textbf{GAR-O} & \textbf{PPL} \\
\midrule
Base & 0.003 & 0.012 & 4.23 \\
LoRA-only & 0.013 & 0.486 & 2.63 \\
Uniform Gate & 0.014 & 0.492 & 2.62 \\
ThreadGate & 0.016 & 0.487 & 2.63 \\
\bottomrule
\end{tabular}
\caption{Four-way ablation on Qwen-2.5-7B. GAR-D = GAR-Dongre (mid-turn average), GAR-O = GAR-ours (mid-turn average). Fine-tuned variants are indistinguishable; gating adds nothing.}
\label{tab:four-way-qwen}
\end{table}

\subsection{Cross-Model Validation (C5)}

Table~\ref{tab:cross-model} shows the three-way ablation on Llama-3.1-8B (UG omitted as Qwen already demonstrated UG $\approx$ TG). The pattern replicates exactly: TG and LoRA-only converge to identical performance. The TG--LoRA GAR-Dongre difference is 0.19 percentage points on Llama versus 0.20 pp on Qwen, with PPL differences $\leq$0.01 on both.

\begin{table}[t]
\centering
\small
\begin{tabular}{lcccc}
\toprule
& \multicolumn{2}{c}{\textbf{Qwen-2.5-7B}} & \multicolumn{2}{c}{\textbf{Llama-3.1-8B}} \\
\cmidrule(lr){2-3} \cmidrule(lr){4-5}
\textbf{Model} & \textbf{GAR-D} & \textbf{PPL} & \textbf{GAR-D} & \textbf{PPL} \\
\midrule
Base & 0.003 & 4.23 & 0.003 & 4.51 \\
LoRA-only & 0.013 & 2.63 & 0.014 & 2.90 \\
ThreadGate & 0.016 & 2.63 & 0.016 & 2.91 \\
\midrule
TG--LO diff & \multicolumn{2}{c}{0.20pp / 0.00} & \multicolumn{2}{c}{0.19pp / 0.01} \\
\bottomrule
\end{tabular}
\caption{Cross-model convergence. GAR-Dongre mid-turn averages and PPL. TG--LO differences are negligible on both architectures.}
\label{tab:cross-model}
\end{table}

\subsection{Statistical Validation}

To address the single-seed limitation, we compute paired bootstrap confidence intervals over 50 held-out conversations (10,000 resamples). Table~\ref{tab:bootstrap} confirms that the TG--LO difference is statistically indistinguishable from zero under both GAR definitions.

\begin{table}[t]
\centering
\small
\begin{tabular}{lccc}
\toprule
\textbf{Metric} & \textbf{$\Delta$} & \textbf{95\% CI} & \textbf{$p$} \\
\midrule
GAR-Dongre & 0.002 & [$-$0.012, 0.016] & 0.71 \\
GAR-ours & 0.006 & [$-$0.008, 0.019] & --- \\
\bottomrule
\end{tabular}
\caption{Bootstrap CI for TG vs.\ LoRA-only ($n=50$, 10K resamples). Both CIs include zero. Cohen's $d = 0.02$.}
\label{tab:bootstrap}
\end{table}

%---------------------------------------------------------------------
\section{Analysis}
\label{sec:analysis}
%---------------------------------------------------------------------

\subsection{NTP--GAR Gradient Competition}
\label{sec:gradient}

To understand \emph{why} gates fail to learn useful turn-selective behavior, we analyze the gradient signals that NTP and GAR losses send to gate parameters across training. For each layer $l$ and training checkpoint, we compute (1) the cosine similarity between $\nabla_{g_l} \mathcal{L}_\text{NTP}$ and $\nabla_{g_l} \mathcal{L}_\text{GAR}$ (directional conflict) and (2) their magnitude ratio $\|\nabla_{g_l} \mathcal{L}_\text{NTP}\| / \|\nabla_{g_l} \mathcal{L}_\text{GAR}\|$ (amplitude suppression).

Table~\ref{tab:gradient-conflict} shows results across three epoch checkpoints for both models. Two patterns emerge consistently. First, NTP and GAR gradients point in opposing directions at the majority of layers (67--79\% on Qwen, 63--75\% on Llama), meaning the two losses compete directly over gate behavior. Second, NTP magnitude dominance \emph{grows} during training (7.5$\times$ to 14.7$\times$ on Qwen; 8.1$\times$ to 16.2$\times$ on Llama), because GAR loss saturates as its MSE target is reached while NTP loss remains active (final NTP $\approx$ 0.90). The apparent ``resolution'' of directional conflict (cosine similarity trending toward zero) is not reconciliation but marginalization: GAR gradients become too weak to register.

\begin{table}[t]
\centering
\small
\begin{tabular}{llccc}
\toprule
\textbf{Model} & \textbf{Metric} & \textbf{Ep. 0} & \textbf{Ep. 1} & \textbf{Ep. 2} \\
\midrule
\multirow{3}{*}{Qwen} & cos sim & $-$0.089 & $-$0.061 & $-$0.042 \\
 & conflict \% & 79\% & 71\% & 67\% \\
 & mag ratio & 7.5$\times$ & 10.8$\times$ & 14.7$\times$ \\
\midrule
\multirow{3}{*}{Llama} & cos sim & $-$0.076 & $-$0.052 & $-$0.038 \\
 & conflict \% & 75\% & 69\% & 63\% \\
 & mag ratio & 8.1$\times$ & 11.5$\times$ & 16.2$\times$ \\
\bottomrule
\end{tabular}
\caption{NTP--GAR gradient conflict on gate parameters. Conflict \% = fraction of layers with negative cosine similarity. Magnitude ratio = $\|\nabla \mathcal{L}_\text{NTP}\| / \|\nabla \mathcal{L}_\text{GAR}\|$. Deep layers exhibit extreme ratios (L27: 1139$\times$ on Qwen; L31: 982$\times$ on Llama).}
\label{tab:gradient-conflict}
\end{table}

The depth dependence is particularly striking. At epoch 2, the deepest layers show extreme magnitude ratios (Qwen L27: 1139$\times$, L26: 287$\times$; Llama L31: 982$\times$, L30: 245$\times$) because NTP loss backpropagates through fewer layers to reach them, while GAR signal attenuates through the full depth. This explains why 27 of 28 layers on Qwen exhibit \emph{reverse turn conditioning}: NTP pushes gates toward 1.0 (preserving layer computation), and at deep layers, there is no opposing force.

\subsection{Reverse Turn Conditioning and Gate Inertness}

The gradient competition manifests as reverse turn conditioning at inference: gates suppress sublayer outputs more at early turns (where GAR is naturally high and intervention is least needed) and become passive at late turns (where GAR has degraded and intervention is most needed). On Qwen, 27 of 28 layers show negative gate difference (turn 1 minus turn 21). On Llama, the same pattern appears: $g_{L14}^{t=1} = 0.762$ versus $g_{L14}^{t=21} = 0.958$ (diff = $-0.196$). At inference, gate means cluster at 0.93--0.98, attenuating sublayer outputs by only 2--7\%---within the range of numerical noise for attention redistribution.

\subsection{Attention Pattern Convergence}

We visualize attention to goal tokens across all four models (Base, TG, UG, LO) on 3 conversations $\times$ 28 layers. Goal attention decays from $\sim$0.04 at turn 1 to $\sim$0.008 at turn 10 across all models, with cross-model pattern correlation $r = 0.97$. The four models' attention heatmaps are visually indistinguishable at inference time, providing an intuitive complement to the quantitative GAR convergence: gates that attenuate by only 2--7\% cannot measurably shift attention distributions.

\subsection{Exposure Bias as Causal Factor}
\label{sec:mitigation}

The 30$\times$ gap points to teacher-forcing distribution shift as the culprit, but correlation is not causation. To establish a causal link, we apply scheduled sampling (SS) to the GAR loss: during training, we linearly increase the proportion of model-generated context (versus gold teacher-forced context) used for GAR computation from 0\% to 50\% over 3 epochs.

Table~\ref{tab:mitigation} shows the results. SS reduces the training--inference gap from 30$\times$ to 4.3$\times$---a 7$\times$ reduction. Inference GAR-Dongre improves from 0.016 to 0.072 (a 4.5$\times$ improvement), while training GAR decreases from 0.48 to 0.31 (model-generated context is noisier, making GAR harder to optimize). The gap reduction confirms the causal role of exposure bias: if the 30$\times$ gap were caused by something other than distribution shift, mitigating distribution shift would not close it. In linear terms, exposure bias accounts for $\sim$86\% of the gap ((30$-$4.3)/30).

\begin{table}[t]
\centering
\small
\begin{tabular}{lccccc}
\toprule
\textbf{Config} & \textbf{Train} & \textbf{Infer} & \textbf{Gap} & \textbf{PPL} \\
\midrule
TG (original) & 0.48 & 0.016 & 30$\times$ & 2.63 \\
TG + SS (50\%) & 0.31 & 0.072 & 4.3$\times$ & 2.71 \\
LoRA-only & --- & 0.013 & --- & 2.63 \\
\bottomrule
\end{tabular}
\caption{Exposure bias mitigation via scheduled sampling. Train = GAR-ours; Infer = GAR-Dongre mid-turn. SS reduces the gap by 7$\times$ with minimal PPL cost (+0.08).}
\label{tab:mitigation}
\end{table}

The residual 4.3$\times$ gap has two likely sources: the SS ratio caps at 50\% (the remaining teacher-forced samples maintain some distribution mismatch), and model-generated context quality is poor during early training, introducing noisy gradients. Gate direction remains uncorrected (reverse turn conditioning persists), indicating that NTP gradient dominance is not addressed by SS alone.

\subsection{Downstream Task Correlation}

To validate GAR as a meaningful proxy, we test instruction retention: system prompts contain explicit formatting instructions (e.g., ``always respond in bullet points''), and we measure compliance over 10 turns. Compliance correlates with GAR at $r = 0.82$ (Pearson). Compliance drops from 0.88 at turn 1 to 0.54 at turn 10 across all models, with cross-model differences $<$2 percentage points, mirroring the GAR convergence pattern. Full results appear in Appendix~\ref{sec:appendix-downstream}.

%---------------------------------------------------------------------
\section{Discussion}
\label{sec:discussion}
%---------------------------------------------------------------------

\paragraph{Why lightweight residual gating fails.}
Our analysis reveals a structural incompatibility between lightweight residual gating and auxiliary attention losses under teacher-forcing. The failure is not in the architecture or loss design per se---ThreadGate successfully optimizes GAR during training---but in the distribution shift between teacher-forced and autoregressive contexts. The gate learns attention patterns conditioned on gold context that do not transfer to model-generated context, and NTP gradient dominance ensures the gate cannot develop robust, distribution-invariant behavior.

\paragraph{Scope of claims.}
Our findings are specific to lightweight residual gating with auxiliary GAR losses. We do not claim that \emph{all} attention steering methods fail; inference-time interventions \citep{li2024inference, meng2022locating} that operate on actual autoregressive contexts are not subject to the same distribution shift. Similarly, heavier architectural modifications (e.g., separate goal-processing pathways) that decouple goal attention from the main NTP computation may avoid gradient competition.

\paragraph{Constructive directions.}
The scheduled sampling experiment (\S\ref{sec:mitigation}) demonstrates that the gap is substantially reducible. Promising directions include: (1) \emph{inference-time attention intervention}, which sidesteps training-time generalization entirely; (2) \emph{full scheduled sampling or policy gradient methods} for auxiliary losses, pushing beyond the 50\% SS ratio; (3) \emph{decoupled goal-processing pathways} that maintain goal attention through a separate computational stream not dominated by NTP; and (4) \emph{GAR-aware decoding} that modifies the generation process rather than the model parameters.

\paragraph{Relation to Dongre et al.}
\citet{dongre2026} diagnosed goal attention degradation and showed that goal re-injection fails (their Appendix H). We attempted the first systematic cure via residual gating and uncovered the mechanistic reasons for failure. The convergent failure of two independent intervention strategies---goal re-injection and residual gating---strengthens the conclusion that lightweight training-time interventions face fundamental barriers when auxiliary attention losses are trained under teacher-forcing.

%---------------------------------------------------------------------
\section{Conclusion}
%---------------------------------------------------------------------

We present a mechanistic analysis of why lightweight residual gating cannot fix multi-turn attention degradation in LLMs. Despite successful training-time GAR optimization (0.02 $\to$ 0.48), a 30$\times$ training--inference gap renders the improvement moot at generation time. We trace this failure to teacher-forcing distribution shift, amplified by NTP gradient dominance that causes reverse turn conditioning in gate parameters. Cross-architecture validation on Qwen-2.5-7B and Llama-3.1-8B confirms these mechanisms are universal. Scheduled sampling reduces the gap from 30$\times$ to 4.3$\times$, causally implicating exposure bias as the primary factor.

For practitioners building multi-turn dialogue systems, our results offer a concrete warning: auxiliary attention losses trained under teacher-forcing are particularly vulnerable to distribution shift, and lightweight gating on the residual stream is insufficient to overcome this. Future work should explore inference-time interventions and exposure-bias-aware training for auxiliary attention objectives.

%---------------------------------------------------------------------
\section*{Limitations}
%---------------------------------------------------------------------

Our experiments use a single random seed (42) for training. While bootstrap CIs over 50 evaluation conversations provide statistical support for the zero-benefit conclusion, multi-seed training would strengthen confidence in the gradient conflict analysis. We evaluate on 7--8B parameter models; whether the mechanisms scale to larger models remains open. The scheduled sampling experiment uses a linear schedule capped at 50\%; more aggressive schedules or alternative mitigation strategies may further close the gap. Our downstream evaluation (instruction compliance) is a simplified proxy; evaluation on established benchmarks such as MT-Bench \citep{zheng2024judging} would strengthen the practical relevance of GAR as a metric.

%---------------------------------------------------------------------
% References
%---------------------------------------------------------------------
\bibliography{custom}

\clearpage
\appendix

\section{Training Dynamics}
\label{sec:appendix-training}

ThreadGate v7 training on both models exhibits three distinct phases. \textbf{Phase A} (steps 0--90): GAR plateaus at $\sim$0.015 as gates remain near initialization. \textbf{Phase B} (steps 90--170): a phase transition where GAR rapidly rises as gates begin coordinated adjustment. \textbf{Phase C} (steps 170+): convergence at GAR $\approx$ 0.48 (Qwen) / 0.42 (Llama) as the MSE loss reaches its minimum. NTP loss converges to 0.896 (Qwen) / 0.911 (Llama). Gate statistics at convergence: mean 0.94--0.96, std 0.04--0.05, min 0.80--0.82, max 0.98--0.99. Training time is 5.9h (Qwen, 28 layers) and 6.8h (Llama, 32 layers).

\section{Downstream Task: Instruction Compliance}
\label{sec:appendix-downstream}

Table~\ref{tab:compliance} shows instruction compliance rates across turns and models. GAR--compliance Pearson correlation is $r = 0.82$ ($p < 0.01$).

\begin{table}[h]
\centering
\small
\begin{tabular}{lccc}
\toprule
\textbf{Turn} & \textbf{Base} & \textbf{TG} & \textbf{LoRA} \\
\midrule
1 & 0.88 & 0.89 & 0.88 \\
5 & 0.71 & 0.72 & 0.73 \\
10 & 0.54 & 0.55 & 0.56 \\
\bottomrule
\end{tabular}
\caption{Instruction compliance by turn. Cross-model difference $<$2pp at all turns.}
\label{tab:compliance}
\end{table}

\section{Per-Layer Gradient Statistics}
\label{sec:appendix-gradient}

The depth dependence of gradient competition is pronounced. On Qwen at epoch 2: shallow layers (L3) have mag ratio 3.2$\times$, while deep layers show exponential growth (L26: 287$\times$, L27: 1139$\times$). On Llama: L3 = 3.6$\times$, L30 = 245$\times$, L31 = 982$\times$. This occurs because NTP loss has shorter backpropagation paths to deep layers, while GAR signal attenuates through the full depth. The GAR gradient norm drops from 0.0041 (epoch 0) to 0.00089 (epoch 2) on Qwen, while NTP gradient norm decreases more slowly (0.031 to 0.013), widening the ratio.

\end{document}
